{\color{red}
\section{Quantum Gate Complexity Analysis}
\label{section:7}

The resource analysis distinguishes the logical encryption operations from
their synthesis into a specified gate library. Let $N=HW$ be the number of
valid primary pixel positions, $n=\lceil\log_2N\rceil$, and $r$ the number
of feature qubits. The image register contains $Q=n+r$ qubits; working
ancillas must be counted separately. The following counts concern the
three encryption layers after image preparation and key generation.
Neither an image-register dimension nor a circuit gate count is an
attack-complexity bound.

\subsection{Register allocation and logical operations}

Table~\ref{tab:1} fixes the register conventions used in the resource and
security comparisons. Here $s$ is the binary intensity depth, $p$ the
adjacency-register size, $c$ the number of channels, and $b$ the number of
wavelength-label qubits. A register with one possible wavelength label has
$b=0$. These explicit benchmark choices do not identify the allocation
of a hardware experiment. The HSI and adjacency allocations follow their
position--feature descriptions
\cite{QIRHSI,EQIRHSI,AFQIRHSI}; the QIRBP and QIRWC layouts distinguish
bit-plane, wavelength, and neighbouring-position registers
\cite{QIRBPQIRWC}. QRMW denotes the multi-wavelength representation with
$s$ intensity qubits and $\lceil\log_2c\rceil$ channel-label qubits
(\url{https://doi.org/10.3906/elk-1705-396}); it is distinct from QIRWC.

\begin{table}[htbp]
\centering
\color{red}
\caption{{\color{red}Declared position--feature allocations for QEF resource accounting.}}
\label{tab:1}
\small
\renewcommand{\arraystretch}{1.15}
\begin{tabular}{|l|c|c|c|}
\hline
\rowcolor{Bisque3}
\textbf{Model} & \textbf{Feature size $r$} & \textbf{Benchmark} & \textbf{$Q$} \\
\hline
FRQI & $1$ & $r=1$ & $n+1$ \\\hline
NEQR & $s$ & $s=8$ & $n+8$ \\\hline
QIRHSI & $s+2$ & $s=8$ & $n+10$ \\\hline
EQIRHSI & $s+2$ & $s=8$ & $n+10$ \\\hline
AFQIRHSI & $p+3$ & $p=4$ & $n+7$ \\\hline
QIRBP & $b+6$ & $b=0$ & $n+6$ \\\hline
QRMW & $s+\lceil\log_2c\rceil$ & $s=8,\ c=3$ & $n+10$ \\\hline
QIRWC & $n+2b+8$ & $b=0$ & $2n+8$ \\\hline
\end{tabular}
\end{table}

For the QIRWC row, the primary position register has $n$ qubits and
the neighbouring-position register is included in the $r$ feature
qubits. A prescribed neighbour $f(k)$ is represented within the
normalised feature state $|C_k\rangle$, so the input retains the form
$N^{-1/2}\sum_k|k\rangle|C_k\rangle$. The row does not use an
independent address register ranging over all $N^2$ pixel pairs. The
QIRBP allocation counts three RGB data qubits, three bit-plane-index
qubits, and $b$ wavelength-label qubits; it does not identify the number
of bit planes with the number of index qubits.

The pixel-dependent substitution specified in
Theorem~\ref{Generalized Quantum Image Encryption via HG-QRNG} contains
three address-controlled rotations for each of the $r$ feature qubits
at each of the $N$ valid positions. Its specified logical count, before
identity removal or circuit optimisation, is therefore
\begin{equation}
L_{\rm sub}=3Nr.
\label{eq:qef-logical-sub-count}
\end{equation}
The controls select $|k\rangle$; they are not ordinary unconditioned
single-qubit gates. Invalid addresses receive the identity operation.
The entanglement loop in Algorithm~\ref{Alg2} selects one CNOT or CZ
for each position--feature qubit pair. Consequently,
\begin{equation}
L_E=L_{CX}+L_{CZ}\le nr,
\qquad L_E=nr\ \text{for the complete loop}.
\label{eq:qef-logical-ent-count}
\end{equation}
This loop is indexed by position qubits and feature qubits, rather than
by $N$ pixel addresses. Its count is therefore $nr$, not $Nr$.

\begin{table}[htbp]
\centering
\color{red}
\caption{{\color{red}Logical operation counts for the allocations in Table~\ref{tab:1}.}}
\label{tab:QGC}
\small
\renewcommand{\arraystretch}{1.12}
\begin{tabular}{|l|c|c|}
\hline
\rowcolor{Bisque3}
\textbf{Model} & \shortstack{\textbf{Controlled rotations}\\$L_{\rm sub}$}
& \shortstack{\textbf{Entangling gates}\\$L_E$ (complete loop)} \\
\hline
FRQI & $3N$ & $n$ \\\hline
NEQR & $24N$ & $8n$ \\\hline
QIRHSI / EQIRHSI & $30N$ & $10n$ \\\hline
AFQIRHSI & $21N$ & $7n$ \\\hline
QIRBP ($b=0$) & $18N$ & $6n$ \\\hline
QRMW ($s=8,c=3$) & $30N$ & $10n$ \\\hline
QIRWC ($b=0$) & $3N(n+8)$ & $n(n+8)$ \\\hline
\end{tabular}
\end{table}

\subsection{Compiled gate count and depth}

Fix a gate library, hardware connectivity, ancilla policy, and rotation
synthesis accuracy. Let $g_C(n)$ and $d_C(n)$ be upper bounds on the
compiled gate count and depth of one address-controlled rotation,
including any control-bit inversions and restoration of its working
ancillas. Their values depend on the chosen synthesis; no specific
asymptotic formula follows from the abstract operator definition alone.
Let $g_E$ and $d_E$ bound the compiled resources of one selected CNOT
or CZ. These bounds equal one for count and depth when both gates are
available primitives with the required connectivity.

Denote the compiled gate count and depth of a clean implementation of
$P_\pi$ by $G_P$ and $D_P$. These quantities include the chosen key
lookup, comparator or routing construction, its scratch registers,
and their cleanup whenever those operations are part of the
permutation implementation. In particular, a sorting description does
not by itself give an $N$-tuple sorting-network cost for a single
amplitude-encoded image register. A construction must specify how the
keyed bijection and its inverse are accessed and implemented. The
permutation must realise
\[
|k\rangle|0\rangle_{\rm work}
\longmapsto
|\pi(k)\rangle|0\rangle_{\rm work};
\]
retaining address-dependent garbage would change the encryption
channel when that garbage is inaccessible to the recipient.

For the concatenation of the three compiled modules, serial execution
of the controlled rotations and selected entangling gates gives the
constructive upper bounds
\begin{align}
\mathcal{QGC}_{\rm QEF}
&\le 3Nr\,g_C(n)+G_P+L_Eg_E,
\label{eq:qef-compiled-gate-bound}\\
\mathcal{QCD}_{\rm QEF}
&\le 3Nr\,d_C(n)+D_P+L_Ed_E.
\label{eq:qef-compiled-depth-bound}
\end{align}
Parallel scheduling or cross-module optimisation may improve these
bounds. They are not asserted to be optimal. The serial-module depth
is the sum of the three scheduled module depths, while its gate count
is the sum of their gate counts.

For a primitive type $\chi$ in the chosen library $\mathcal L$, let
$g_\chi(k,j,a)$ count its occurrences in the synthesis of the
address-controlled $R_a(\theta_{k,j}^a)$ gate. The precise accounting
identity is
\begin{equation}
G_\chi=
\sum_{k=0}^{N-1}\sum_{j=1}^{r}\sum_{a\in\{x,y,z\}}
g_\chi(k,j,a)+G_\chi^{(P)}+G_\chi^{(E)}.
\label{eq:84}
\end{equation}
In particular, CNOT gates used to synthesise controlled substitution
rotations belong to the substitution block, not the permutation block.
For a library containing CNOT and CZ directly, the unsynthesised
entanglement counts are $G_{CX}^{(E)}=L_{CX}$ and
$G_{CZ}^{(E)}=L_{CZ}$; otherwise their chosen decompositions must be
included in Eq.~\eqref{eq:84}.

\subsection{Weighted circuit cost}

\begin{theorem}[Additive cost of a specified QEF implementation]
\label{thm:QEF-quantum-cost}
For any compiled QEF circuit over a fixed library $\mathcal L$, assign
a nonnegative cost $w_\chi$ to each primitive type $\chi$. With the
counts in Eq.~\eqref{eq:84}, its weighted cost is
\begin{equation}
QC_{\rm QEF}=\sum_{\chi\in\mathcal L}w_\chi G_\chi.
\label{eq:QC-QEF-final}
\end{equation}
For a concatenation retaining its three module boundaries,
\begin{equation}
QC_{\rm QEF}=QC_{\rm sub}+QC_{\rm perm}+QC_{\rm ent}.
\label{eq:QC-sum-blocks}
\end{equation}
\end{theorem}

\begin{proof}
For each module $B\in\{{\rm sub},{\rm perm},{\rm ent}\}$, define
\begin{equation}
QC_B=\sum_{\chi\in\mathcal L}w_\chi G_\chi^{(B)}.
\label{eq:QC-sum-def}
\end{equation}
Each gate occurrence belongs to exactly one module, so
$G_\chi=\sum_B G_\chi^{(B)}$. Substitution and interchange of the
finite sums give
\begin{equation}
\sum_{\chi\in\mathcal L}w_\chi G_\chi
=\sum_B\sum_{\chi\in\mathcal L}w_\chi G_\chi^{(B)}
=\sum_B QC_B,
\label{eq:QC-expanded}
\end{equation}
which proves both identities.
\end{proof}

Unit weights recover the compiled gate count. Hardware durations,
two-qubit-gate weights, or another declared resource model give
different costs. In particular, neither CZ nor Toffoli has a universal
fixed optimal cost independent of the primitive library. Logical
multi-qubit gates and their decomposed primitives must not both be
counted for the same operation.

For an end-to-end circuit, image preparation, generation and loading of
the realised keys, and any required final measurement must be added
separately. When $N<2^n$, the exact amplification construction uses one
phase-matched amplification step, but synthesis of its reflections and
of the feature-loading operations still has size-dependent cost.
The intensity depth $s$ remains independent of the number of pixels
$N$; setting $s=N$ would describe a changing encoding rather than a
fixed eight-bit image model. Numerical preparation estimates obtained
under that identification are therefore not QEF encryption costs.
Numerical compiled counts and depths require a specified clean
permutation circuit and controlled-rotation synthesis; the logical
counts in Table~\ref{tab:QGC} are the quantities established here.
\par}
